\section{Exactness, Refinements and Over-approximation}\label{sec:s-exactness}

The decision $D$ analysed here is the one the implementation runs as ``claim pattern with local
filters'': for each committed version $v\in(r_0,r]$ it computes the restricted triggers of
Corollary~\ref{cor:filters} with respect to $\G_{v-1}$ and $\G_v$, and marks the view if one
exists, at answer level if one of them is an answer trigger and at citation level otherwise. By
Corollaries~\ref{cor:stepwise} and~\ref{cor:filters}, $D$ has no false negatives. The question is
when it has no false positives. The level of $D$ over the interval is the strongest level of
the marks raised in it.

\paragraph{Conditions on the query.}
\begin{description}
\item[(Q1)] One positive pattern $P$, no negated pattern, a projection at the root (counting
  and selection are treated in Remarks~\ref{rem:count} and~\ref{rem:latest}), and provenance
  expansion on (without provenance every $\ev_t$ is empty and every (T4) trigger over-marks).
\item[(Q2)] Current state (not historical) and no world time.
\item[(Q3)] No variable occurs twice in $P$ (the speaker variable included), and the scope
  position is not a variable.
\item[(Q4)] No filter uses $\mathrm{end}(\lambda_P)$, and the matcher decides every filter on the
  written fields exactly as the model does. The second part is partly a condition on the data:
  no claim reached by a trigger has an unknown start compared by a $\mathrm{start}$ filter, and
  no date is compared with a date-time.
\end{description}
Under (Q3) the matcher's relation coincides with $\matches$; under (Q4) the value of a filter on a
row does not depend on the version, since it reads only fields fixed when entries are written.

\paragraph{Conditions on the change.}
\begin{description}
\item[(D1)] No claim reached by a (T1) or (T2) trigger is superseded in $\G_r$.
\item[(D2)] No (T3) trigger points to a claim already superseded in $\G_{r_0}$.
\item[(D3)] No (T4) trigger points to a claim superseded in $\G_r$, and no (T4) trigger on a claim
  $c$ has an identifier already in $\ev_t(\cut{\Q}{r_0},\G_{r_0})$ for an answer $t$ whose witnesses
  contain $c$.
\end{description}

\begin{proposition}[Exactness; Proposition~2 of the main text]\label{prop:s-exactness}
Assume A1--A9, (Q1)--(Q4) and (D1)--(D3). Then (a)~$D$ raises an answer-level mark iff the
answers or the witnesses change between $r_0$ and $r$; (b)~$D$ marks the view at all iff the
answers, the witnesses or the evidence change; in particular, if the answers and witnesses are
unchanged and the evidence changes, $D$ marks the view and no mark it raises in $(r_0,r]$ is at
answer level (the strongest mark is at citation level).
\end{proposition}
\begin{proof}
The ``if'' directions are soundness (Corollaries~\ref{cor:stepwise} and~\ref{cor:filters}). For
the ``only if'' directions, note first that under (Q1) every witness is $\{c\}$ or $\{c,a\}$ for a
row $(\theta,c,a)$, and two distinct such sets are incomparable, so every row's witness is a
minimal witness of its answer.

(a) Let $e$ be a restricted answer trigger at version $v$.
\emph{(T1):} $e=c'$ is a new claim and the row $(\theta,c',\ast)$ satisfies (m1), (m2) and
$\Phi_P$ in $\G_v$, hence in $\G_r$ (Q4). By (D1), $\sups_{\G_r}(c')=\emptyset$, so (m3) holds, and
(m4) is void (Q2). The row is a body row in $\G_r$, so $\{c'\}$ is a witness at $r$; since $c'\notin
\G_{r_0}$ it is no witness at $r_0$.
\emph{(T2):} $e=a'$ is a new attribution on $c$ and $(\theta,c,a')$ satisfies (m1), (m2) and
$\Phi_P$; by (D1) $c$ is current in $\G_r$, so $\{c,a'\}$ is a witness at $r$ and, $a'$ being new,
not at $r_0$.
\emph{(T3):} $e$ is a supersession event on $c$, and some row $\mu=(\theta,c,a)$, ignoring (m3)
and (m4), satisfies $\Phi_P$. If $c$ were new, its own (T1) trigger, or the (T2) trigger of the
attribution in $\mu$, would reach a claim superseded in $\G_r$, against (D1); so $c\in\G_{r_0}$. If
$a\neq\ast$ were new, it would be a (T2) trigger reaching $c$, again against (D1); so
$a\in\G_{r_0}$. By (D2), $c$ is current in $\G_{r_0}$, so $\mu$ is a body row at $r_0$ and
$W(\mu)\ni c$ is a witness at $r_0$. In $\G_r$, $c$ is superseded, so no row and no witness at $r$
contains $c$.
In every case the witnesses change.

(b) If some trigger is an answer trigger, apply (a). Otherwise there is a (T4) link $\ell$ on a
claim $c$ with a row $(\theta,c,a)$ satisfying (m1), (m2) and $\Phi_P$ in $\G_v$. By (D3), $c$ is
current in $\G_r$, so the row is a body row at $r$ and $c$ is a witness claim; hence
$\mathrm{id}(\ell)\in\ev_t(\cut{\Q}{r},\G_r)$ for the answer $t$ of that row. With no answer
trigger the answers and witnesses are unchanged (soundness), so $t$ is an answer at $r_0$ with $c$
in its witnesses, and by (D3) $\mathrm{id}(\ell)\notin\ev_t(\cut{\Q}{r_0},\G_{r_0})$; the view is
citation-stale.
The last claim of the statement is the contrapositive of the ``only if'' part of (a): if answers
and witnesses are unchanged, no answer-level mark is raised, and by the ``if'' part of (b) some
mark is.
\end{proof}

\noindent The level of a mark is not exact even under these conditions: an answer-level mark
can come from a change of witnesses alone (Example~\ref{ex:level}). Deciding the level needs the
answer fingerprint, which the hybrid strategy computes after the mark.

\begin{remark}[Distinct count]\label{rem:count}
Proposition~\ref{prop:s-exactness} also holds when the root of $\Q$ is a distinct count, grouped
or not, with (Q1)--(Q4) and (D1)--(D3) otherwise unchanged.
\end{remark}
\begin{proof}
In the cases (T1) and (T2) of (a), the new row adds $c'$ (resp.\ $a'$) to the union witness of its
group, or creates the group; in case (T3), $c$ lies in the union witness of its group at $r_0$ and
in no row at $r$, so it leaves the union or the group disappears. Part (b) is unchanged.
\end{proof}

\begin{remark}[Latest by version]\label{rem:latest}
Under (Q2)--(Q4), for a latest-state selection by version over the single pattern, every
restricted (T1) or (T2) trigger (a claim or attribution insertion) satisfying (D1) changes the
witnesses. A (T3) trigger on a claim that is not selected at $r_0$ need not
(Example~\ref{ex:latest-version}), nor need a (T4) trigger on an unselected claim, whose evidence
is not cited.
\end{remark}
\begin{proof}
The new row $\mu'$ found in (a) contains an entry $e\in\Dlt$, so its key is at least
$\vs(e)>r_0$, while every key at $r_0$ is at most $r_0$. At $r$ the maximum key over the rows is
therefore greater than $r_0$, and each selected claim has a kept row whose key comes from an entry
of $\Dlt$; the witness of that row contains an entry of $\Dlt$ and so is no witness at $r_0$.
\end{proof}

\subsection{Over-approximation outside the conditions}

Each example gives a query, a graph at $r_0$ and a change $\Dlt$ for which $D$ marks the view
while answers, witnesses and evidence are unchanged (unless stated otherwise). Entity and predicate
names follow the running examples of the main text.

\begin{example}[Join without a partner]\label{ex:join}
$\Q$: \texttt{SELECT ?t ?arr WHERE \{ (req1, selectedService, ?t)[user, accept] .
(?t, scheduledArrival, ?arr) \}}. $\Dlt$ adds the claim \texttt{(TR9, scheduledArrival,
"21:00:00"\textasciicircum\textasciicircum time)} with its attribution and link, while TR9 was
never selected. The second pattern has no bracket, so the claim is a (T1) trigger of it, but joins with no row of the first.
\end{example}

\begin{example}[Negated pattern without an answer]\label{ex:neg}
$\Q$: \texttt{SELECT ?t WHERE \{ (req1, selectedService, ?t)[system, offer] NOT EXISTS \{
(req1, selectedService, ?t)[user, accept] \} \}}. $\Dlt$ adds a claim
\texttt{(req1, selectedService, TR9)} stated and accepted by the user, never offered. It triggers
the negated pattern, but no positive row has $t=\mathrm{TR9}$.
\end{example}

\begin{example}[Latest by valid time]\label{ex:latest-valid}
$\Q$: \texttt{SELECT ?city WHERE \{ (lina, livesIn, ?city) AS ?c \} LATEST BY valid OF ?c}, with
the selected claim starting in 2023. $\Dlt$ adds a current claim of the pattern starting in 2019:
a (T1) trigger, but it is not the latest.
\end{example}

\begin{example}[Latest by version, supersession of an unselected claim]\label{ex:latest-version}
Same pattern selected by version, with two current claims $c_1$ (stamp 1) and $c_2$ (stamp 2);
$c_2$ is selected. $\Dlt$ supersedes $c_1$ (a removal on a multi-valued predicate): a (T3)
trigger, while the rows of $c_2$ are unchanged.
\end{example}

\begin{example}[Historical state]\label{ex:hist}
$\Q$: \texttt{SELECT ?city WHERE \{ (lina, livesIn, ?city) \} AT(state=historical)}. $\Dlt$ adds a
change event on a matching claim: a (T3) trigger, but (m3) excludes only corrected claims, so the
rows are unchanged. This holds only without world time and without $\mathrm{end}$ filters: with
\texttt{AT(valid="2022-06-01"\textasciicircum\textasciicircum date, state=historical)}, a change
event stating $\vend=$ 2022-03-01 on a claim with open $\tau^+$ sets its effective end before $d$,
so (m4) removes the claim and the answers \emph{do} change.
\end{example}

\begin{example}[Repeated supersession; superseded within the change]\label{ex:resup}
(i)~A claim superseded at $r_0$ receives a second supersession event in $\Dlt$ (violates (D2)).
(ii)~A new matching claim and its supersession both lie in $\Dlt$ (violates (D1)). (iii)~An
acceptance is added to a claim that is already superseded (violates (D1)).
\end{example}

\begin{example}[Evidence]\label{ex:evidence}
(i)~A new link on a matching claim cites a span whose identifier is already cited for the same
answer, for example the same message cited again for the same claim; $\ev_t$ is a set of
identifiers and does not change (violates (D3)). (ii)~A new link on a superseded matching claim (violates (D3)).
\end{example}

\begin{example}[Level of a mark]\label{ex:level}
$\Q$: \texttt{SELECT ?t WHERE \{ (req1, selectedService, ?t)[user, accept] \}} with answer TR7143.
$\Dlt$ adds a second (user, accept) attribution on the same current claim. It is a (T2) trigger and
$D$ raises an answer-level mark; the answers are unchanged and the witnesses gain
$\{c,a'\}$, so the view is only citation-stale.
\end{example}

\noindent Conditions (Q3) and (Q4) exclude the remaining sources: a repeated variable or a scope
variable (the matcher ignores both), an undecidable filter (counted as satisfied), and an
$\mathrm{end}$ filter on a supersession trigger (tried before and after the change, although a
current-state query drops the superseded claim anyway).

\subsection{The observation}

Let $U$ be the policy that marks a view iff $\Dlt$ contains a supersession event on a claim of its
witnesses at $r_0$, and $U^+$ the policy that also marks it on a trigger of its negated pattern or
of a pattern under a count.

\begin{proof}[of the Observation in the main text]
Each case is an instance with no supersession event in $\Dlt$, so $U$ does not mark.
(i)~$\Q$: \texttt{SELECT ?s WHERE \{ (lina, responsibleFor, ?s) \}}; $\G_{r_0}$ holds the claim for
Atlas; $\Dlt$ adds the claim for Boreal with an own positive attribution and a link. The answers
gain Boreal.
(ii)~$\Q$ of Example~\ref{ex:neg}; $\G_{r_0}$ holds the claim for TR7143 with a (system, offer)
attribution, so the answer is TR7143; $\Dlt$ adds a (user, accept) attribution on that claim. The
answer set becomes empty.
(iii)~$\Q$: \texttt{SELECT COUNT(DISTINCT ?i) WHERE \{ (melanie, playsInstrument, ?i) \}};
$\Dlt$ adds the clarinet claim to the violin claim. The count changes from 1 to 2.
$U^+$ marks (ii) and (iii) but not (i), whose only trigger is a (T1) trigger of a positive pattern
without counting. In each case the missed entry is an insertion that supersedes nothing: a new
claim, or in (ii) a new attribution on a used claim. A policy that also watched attributions and
links on used claims would catch (ii), but still not (i) or (iii).
\end{proof}

\subsection{What is not claimed}

The exact decision for a given graph and change is re-evaluation: read-time re-query compares the
answers, witnesses and evidence directly and is exact with respect to the claim graph. We make no
claim about the decidability or complexity of \emph{static} query--update independence for this
fragment; the undecidability results of Levy and Sagiv~\cite{levy1993independent} concern more
expressive query languages. Nor do we claim a new invalidation algorithm, the first soundness
result for SPARQL cache invalidation, or an incremental maintenance algorithm. For monotone queries, deleting an
entry outside every witness leaves the answers unchanged, which follows from why-provenance
\cite{buneman2001why}, so a citation-based policy is correct for them under deletion; an append-only graph has no deletion, supersession events
play its role through the built-in negation of every current-state query, and they are exactly
the insertions such a policy observes.
